Le stage réalisé au sein de l'Office National des Pêches avait pour objectif de concevoir et de réaliser une application web de gestion des ventes de blocs de glace, en réponse à une gestion jusqu'alors manuelle, source d'erreurs et dépourvue de traçabilité.

L'application \textbf{IceVentes} atteint cet objectif. Elle centralise l'ensemble du processus : authentification et gestion des utilisateurs selon trois rôles, paramétrage du référentiel et de son prix, gestion des demandeurs, saisie des ventes avec génération automatique d'un code unique et calcul automatisé entre prix et quantité, historique multicritère avec exports, et tableau de bord de pilotage. Une attention constante a été portée à la cohérence des données, à la sécurité des accès et à la qualité de l'expérience utilisateur.

Sur le plan personnel, ce projet a été l'occasion de mettre en pratique et d'approfondir de nombreuses compétences : la conception et la modélisation UML, le développement web avec Django, la manipulation d'une base de données relationnelle, la construction d'interfaces responsives et interactives, ainsi que la prise en compte des exigences de sécurité et de fiabilité propres à une application de gestion. Le travail dans un cadre professionnel a également renforcé la rigueur, l'autonomie et la capacité à structurer un projet de bout en bout.

Plusieurs \textbf{perspectives} d'évolution peuvent être envisagées :

\begin{itemize}
  \item la mise en production sur une base PostgreSQL et un serveur dédié ;
  \item la prise en charge de plusieurs usines à glace actives simultanément ;
  \item l'enrichissement du reporting (statistiques avancées, comparaisons, export planifié) ;
  \item l'ajout d'une gestion fine des paiements et des encaissements ;
  \item le développement d'une version mobile pour la saisie sur le terrain.
\end{itemize}

En définitive, IceVentes constitue une base solide, fonctionnelle et extensible, qui répond aux besoins exprimés tout en ouvrant la voie à des améliorations futures.
